\section{Matched predictive evaluation}\label{sec:evaluation}
\subsection{Finite decoders}
For each encoding $\Gamma$, fit its column means and population standard deviations on A. Scales below $10^{-8}$ are replaced by one; no column is deleted. Let $u(z)\in\R^p$ denote the resulting standardized encoding. Membership-prefix transformations are identical when an enrichment is appended at the same reconstruction and cap. There is no radial bounding map and no forecast clipping.

The linear decoder uses $f(z)=u(z)$. The nonlinear decoder retains this input and appends a Gaussian Nystr\"om block. Choose 64 A landmark row identities, shared across all encodings and caps. Their encoded values $\ell_i\in\R^p$ differ between representations. For bandwidth $\sigma>0$, let $k_\sigma(u,v)=\exp(-\|u-v\|_2^2/(2\sigma^2))$. If the landmark Gram matrix has retained eigenvectors $V$ and positive eigenvalues in the diagonal matrix $\Lambda$, let $k_L(u)$ be the column vector of its 64 landmark similarities. The nonlinear block is $\nu(u)=\Lambda^{-1/2}V^\top k_L(u)$. After its own A-only column standardization, denoted by $\widetilde\nu$, the decoder input is $f(z)=\col(u(z),\widetilde\nu(u(z)))$. The supplement fixes the bandwidth scale and eigenvalue thresholds. Equal landmark counts do not equate effective rank, feature width, or computation.

Let $n_A$ be the number of retained A rows, let $H=|\mathcal H|$, let $Y_i\in\R^H$ contain the A targets, and let $p_f$ be the width of $f$. Fit $a\in\R^H$ and $W\in\R^{p_f\times H}$ by
\begin{equation}
 (\widehat a,\widehat W)=\arg\min_{a,W}
 \left\{\frac1{n_A}\sum_{i\in A}\|Y_i-a-W^\top f(z_i)\|_2^2+\alpha\|W\|_F^2\right\},\quad\alpha>0.
 \label{eq:ridge}
\end{equation}
The intercept is unpenalized. One penalty, and one bandwidth for the nonlinear decoder, is selected jointly across horizons; the coefficient columns remain horizon-specific. Forecasts are direct evaluations of $\widehat a+\widehat W^\top f(z_t)$, not powers of a common transition matrix.

\subsection{Laws, allocation, and candidate selection}
The study has 64 independent repetitions under each of four numerical observation laws: linear autoregression, nonlinear autoregression, Van der Pol, and Duffing. The oscillators include sensor noise. Supplement Section~S1 provides every equation, initialization, burn-in, and integration convention. The laws were used in development studies; fresh repetitions are not new application domains.

Each task contains independently generated A/B/C/D sets with 64/32/32/128 episodes of length 72. Histories have $d=12$ and $\mathcal H=\{1,4,12\}$. Zero-based origins $11,15,\ldots,59$ give thirteen windows per episode. A/B/C retain 512/256/256 uniformly sampled windows without replacement; D retains all 1,664. All procedures share the task's rows. Windows within an episode overlap and are not separate inferential units.

Two A-fitted $k$-means partitions request 2 and 4 groups, with at most 24 references per group. Native cap candidates are $1,2,4,8,12$. The ridge grid is $\{10^{-6},10^{-5},\ldots,10^2\}$ and the nonlinear bandwidth multipliers are $0.25,0.5,1,2,4$. The affinity parameters are fixed at $\tau=10^{-6}$ and $\omega=16$. Every candidate is fitted on A. B minimizes the mean horizon MSE divided by the corresponding retained-A target population variances, each floored at $10^{-12}$. Exact ties use ascending cap, bandwidth, and penalty order. There is no A+B refit.

In the \emph{primary} view, every complete reconstruction--encoding--decoder procedure selects its own cap and settings. The raw-input control has no cap search. In the \emph{secondary} view, affine-local linear ridge selects one common cap; each procedure then selects its own within-cap decoder settings. This is an asymmetric selector, not a neutral bank choice. The primary question concerns deployable selected procedures; the secondary question conditions on a specified common bank. Neither isolates the causal effect of adding a block with every fitted choice fixed.

The sample size, menu, and analysis were fixed locally before production data generation, with knowledge of development results. All A/B fits and selections were sealed before C/D generation, and all C/D predictions before scoring. C supplies diagnostics only. Source hashes document execution order, not external oversight or cryptographic target isolation. No candidate, task, or comparison was excluded. The study is not described as publicly preregistered.

\subsection{Loss, estimand, and the complete primary family}
For law $\ell$, repetition $r$, procedure $m$, and horizon $h$, let $v_{\ell rh}$ be the retained-A target variance with population divisor and floor $10^{-12}$. Let $y_{ih}$ and $\widehat y_{mih}$ be the D target and forecast. Define standardized mean squared error (SMSE) by
\begin{equation}
 L_{\ell rm}=\frac13\sum_{h\in\mathcal H}
 \frac{1}{1664}\sum_{i\in D}\frac{(y_{ih}-\widehat y_{mih})^2}{v_{\ell rh}}.
 \label{eq:loss}
\end{equation}
The estimand is expected terminal loss over the stated finite training, validation, reference-selection, and evaluation procedure for one law and decoder. The random A denominator is part of that procedure. This is not population-optimal decoder risk or a ratio of population expected errors.

Five comparisons are made for every law and decoder: KAHM-local against KAHM-memberships, KAHM-normalizer against KAHM-memberships, and KAHM-local against affine-local, positive-kernel-local, and raw input. These $5\times4\times2=40$ comparisons form one two-sided confidence family. For a candidate $a$ and reference $b$, put $d_r=L_{\ell ra}-L_{\ell rb}$, $\bar d=64^{-1}\sum_r d_r$, and $s_d^2=63^{-1}\sum_r(d_r-\bar d)^2$. The reported interval is
\begin{equation}
 \bar d\ \pm\ t_{63,\,1-0.05/(2\cdot40)}\frac{s_d}{8},
 \qquad t_{63,0.999375}=3.37962646388,
 \label{eq:interval}
\end{equation}
where $t_{\nu,q}$ is the Student quantile with $\nu$ degrees of freedom and probability $q$. A wholly negative interval identifies lower candidate mean loss; a wholly positive interval identifies higher loss. Otherwise the direction is unresolved. A zero sample variance receives no Student interval rather than automatic success. No such exception changes the family size.

Bonferroni adjustment targets simultaneous 95\% coverage when the marginal intervals are valid. Their validity relies on a paired-Student approximation for independent repetition differences and is not distribution-free. No practical, noninferiority, or equivalence margin is specified. Relative gain is $100(1-\overline L_a/\overline L_b)$, a ratio of mean losses, not a confidence endpoint or an average of repetitionwise percentages. Laws are not pooled into a primary aggregate. The 64-repetition budget was the first of $32,48,64,96$ giving a simultaneous half-width below $0.45s_d$; it provides no guarantee of useful absolute precision or resolution of small effects.
